\clearpage
\SourceLeaf{189}
\pdfbookmark[1]{Aufbau und Gebrauch der Tafeln.}{witwen-einrichtung}
\begin{center}AUFBAU UND GEBRAUCH DER TAFELN.\end{center}
\noindent\hspace*{1.2em} [Den Zahlenangaben dieser Tafeln liegen die von \textsc{Brune} im 16. Band des \textsc{Crelle}schen Journals für Mathematik zusammengestellten Beobachtungen zu den von 1776 bis 1834 nacheinander in die Königlich Preußische Allgemeine Witwenverpflegungsanstalt aufgenommenen 31500 Ehepaaren zugrunde. Angegeben sind hier die Logarithmen der Anzahl der Frauen (\(\log fm\)) und der Männer (\(\log FM\)), die von 10000 Personen, die das 20. Lebensjahr vollendet hatten, bis zum Ende des in der Mitte angegebenen Altersjahres (\(m\) oder \(M\)) gelangten. Dabei wird jedoch insofern von \textsc{Brune} abgewichen, als das Sterben der Männer über 80 Jahren nach demselben Verhältnis berechnet ist, das nach jenen Beobachtungen für das weibliche Geschlecht gilt; denn, wie in der Bilanzrechnung von 1845 erwähnt, enthält die Registratur der Preußischen Witwenkasse keine ausreichenden Daten, um das Sterben der Männer im hohen Alter unmittelbar zu bestimmen. Neben den Logarithmen der Lebenden stehen unter der Überschrift decr. die absoluten Werte der Unterschiede dieser Logarithmen (\(\log gm=\log\frac{f(m-1)}{fm}\) und \(\log Gm=\log\frac{F(M-1)}{FM}\)), ausgedrückt in Einheiten der siebten Dezimalstelle. Mithilfe der so erhaltenen Sterblichkeitstafel hat \textsc{Gauss} die einfachen Leibrenten und die Verbindungsrenten bei einem Zinssatz von \(3\tfrac{1}{2}\) Prozent und 4 Prozent berechnet.

Die sowohl in Logarithmen als auch in Zahlen (\(\varphi m,\Phi M\)) angegebenen Werte der einfachen Leibrenten beziehen sich auf das Ende des in der Mitte angegebenen Lebensjahres der Frau (\(m\)) oder des Mannes (\(M\)) als gegenwärtigen Zeitpunkt. Vorausgesetzt wird, dass am Ende jedes folgenden Jahres die Geldeinheit gezahlt wird, sofern die betreffende Person diesen Zeitpunkt erlebt, sodass also
\[
\begin{aligned}
\varphi m\mathord{.}fm&=\rho f(m+1)+\rho\rho f(m+2)+\rho^3f(m+3)+.\ .\ .\\
\Phi M\mathord{.}FM&=\rho F(M+1)+\rho\rho F(M+2)+\rho^3F(M+3)+.\ .\ .
\end{aligned}
\]
gilt, wenn \(\rho\) den Abzinsungsfaktor (\(=\tfrac{200}{207}\) bei \(3\tfrac{1}{2}\) Prozent und \(=\tfrac{25}{26}\) bei 4 Prozent) bezeichnet.

Die Tafel der Verbindungsrenten enthält die Logarithmen der Werte \(\psi(m,M)\). Diese entsprechen für den Zeitpunkt des seitlich angegebenen Alters (\(M\)) des Mannes den Werten solcher Renten, bei denen am Ende jedes folgenden Jahres die Geldeinheit gezahlt wird, sofern der Mann (gegenwärtiges Alter \(=M\)) und die Frau (Altersunterschied \(=m-M\)) gleichzeitig leben. Sie sind also durch die folgende Formel bestimmt:
\[
\begin{aligned}
\psi(m,M)\mathord{.}fm\mathord{.}Fm&=\rho f(m+1)\mathord{.}F(M+1)\\
&\quad+\rho\rho f(m+2)\mathord{.}F(M+2)\\
&\quad+\rho^3f(m+3)\mathord{.}F(M+3)+.\ .
\end{aligned}
\]
\SourceLeaf{190}
\Needspace{5\baselineskip}
\pdfbookmark[2]{Werte von Leibrenten und Lebensversicherungen für Männer.}{witwen-maenner}
\begin{center}\textit{Werte von Leibrenten und Lebensversicherungen für Männer.}\end{center}
\noindent\hspace*{1.2em} Die Tafel für die Leibrenten der Männer hat \textsc{Gauss} dazu benutzt, den Einfluss genau zu berechnen, den die Bestimmung der Statuten hätte, nach der ein lebendes Mitglied nicht wieder austreten dürfte. Er findet, dass sich dadurch für die 42 verheirateten Mitglieder der Bilanzrechnung vom 1. Okt. 1845 der Zeitwert der Beiträge um 898 Taler bei \(3\tfrac{1}{2}\) Prozent und um 665 Taler bei 4 Prozent erhöhen würde. Außerdem hat er mithilfe dieser Tafel einige Berechnungen der Werte von Lebensversicherungen ausgeführt und dabei für das gegenwärtige Alter eines Mannes von \(M\) Jahren \(\rho-(1-\rho)\Phi M\) als Zeitwert der am Ende des Todesjahres auszuzahlenden Geldeinheit angesetzt.]

\Needspace{6\baselineskip}
\pdfbookmark[2]{Wert der bestehenden Witwenpension Pm.}{witwen-pension-pm}
\begin{center}\textit{Wert der bestehenden Witwenpension \(Pm\).}\end{center}
\noindent \(m\)\quad gegenwärtiges Alter der Witwe.\par
\noindent \(\varphi m\)\quad Wert der Pension bei jährlicher Zahlung nur an noch Lebende.\par
\noindent \(\rho\)\quad Abzinsungsfaktor (\(=\tfrac{25}{26}\) für 4 Prozent, \(=\tfrac{200}{207}\) für \(3\tfrac{1}{2}\) Prozent), oder der Zinssatz in dem Sinne verstanden, dass \(\rho\) nach einem Jahr auf 1 anwächst.\par
\noindent \(fx\)\quad Lebende des Alters \(x\) nach den Angaben der Sterblichkeitstafel.

[Die Bilanz wird für den 1. Oktober eines bestimmten Jahres berechnet; dieser Zeitpunkt wird hier überall kurz der gegenwärtige genannt. Nach dem Regulativ vom 11. Oktober 1833 und den späteren Verfügungen zur Professorenwitwenkasse wird die Pension in halbjährlichen, am 1. April und am 1. Oktober jedes Jahres fälligen Raten ausgezahlt. Bei Witwen erlischt sie mit dem Sterbemonat, der noch voll bezahlt wird. Unter Berücksichtigung dessen gilt für den wahrscheinlichen gegenwärtigen Wert \(Pm\) einer Witwenpension, die jährlich die Geldeinheit auszahlt:]
\[
\begin{aligned}
fm\mathord{.}Pm&=\tfrac{1}{12}\rho^{\frac{1}{2}}\{fm+f(m+\tfrac{1}{12})+f(m+\tfrac{2}{12})+.\ .+f(m+\tfrac{5}{12})\}\\
&\quad+\tfrac{1}{12}\rho\{f(m+\tfrac{6}{12})+f(m+\tfrac{7}{12})+.\ .\ .\ .+f(m+\tfrac{11}{12})\}\\
&\quad+\text{usw.}
\end{aligned}
\]
\noindent oder
\[
fm\mathord{.}Pm=\tfrac{1}{2}\rho^{\frac{1}{2}}f(m+\tfrac{5}{24})+\tfrac{1}{2}\rho f(m+\tfrac{17}{24})+\tfrac{1}{2}\rho^{\frac{3}{2}}f(m+\tfrac{29}{24})+\text{usw.}
\]
\noindent wofür man ansetzen kann
\[
=\rho^{\frac{3}{4}}f(m+\tfrac{11}{24})+\rho^{\frac{7}{4}}f(m+\tfrac{35}{24})+\text{usw.}
\]
\noindent Die Tafel gibt
\[
fm\mathord{.}\varphi m=\rho f(m+1)+\rho\rho f(m+2)+.\ .\ .
\]
\noindent also
\[
\begin{aligned}
Pm&=\rho^{-\frac{1}{4}}\mathord{.}\frac{f(m-\tfrac{13}{24})}{fm}\mathord{.}\varphi(m-\tfrac{13}{24})\\
&=\rho^{-\frac{1}{4}}\mathord{.}g(m+\tfrac{11}{48})^{\frac{13}{24}}\mathord{.}\varphi(m-\tfrac{13}{24})
\end{aligned}
\]
\relax
[wenn man die Werte von \(\log f(m-1)-\log fm\), die in der obigen Tafel unter decr. in Einheiten der siebten Dezimalstelle stehen, mit \(\log\mathord{.}gm\) bezeichnet. Eine für die Berechnung von Hilfstafeln etwas bequemere Formel erhält man, wenn man für jede ganze Zahl \(n\) und jeden echten Bruch \(k\) anstelle von \(f(n+k)\) die Größe \((k-1)fn+kf(n+1)\) setzt, nämlich:]
\SourceLeaf{191}
\[
Pm=A+B\mathord{.}\varphi m
\]
\noindent wobei
\[
48A=19\rho^{\frac{1}{2}}+7\rho,\qquad48B=5\rho^{-\frac{1}{2}}+17+19\rho^{\frac{1}{2}}+7\rho
\]
\noindent und hinreichend genau
\[
A=\tfrac{13}{24}\rho^{\frac{33}{52}},\qquad B=\rho^{\frac{7}{24}}
\]
\noindent für \(3\tfrac{1}{2}\) Proz.\quad \(\log\tfrac{1}{\rho}=0.0149403\)\quad Genau
\[
A=0.52998474,\qquad\log B=9.9956901
\]
\noindent die Näherungsformel ergibt
\[
A=0.5299694,\qquad\log B=9.9956424
\]
\noindent Gerechnet worden war nach \(\tfrac{13}{24}+\varphi m=P^0\)\par
\noindent man kann also setzen
\[
\begin{aligned}
P&=\rho^{\frac{7}{24}}P^0-\tfrac{13}{24}(\rho^{\frac{7}{24}}-\rho^{\frac{3}{4}})\\
&=P^0-\tfrac{7}{24}(1-\rho)P^0-\tfrac{143}{576}(1-\rho)
\end{aligned}
\]
\noindent also für \(3\tfrac{1}{2}\) Proz.\quad \(P=P^0-\tfrac{49}{4968}P^0-\tfrac{1091}{119232}\)\par
\noindent für 4 Proz.\quad \(P=P^0-\tfrac{7}{624}P^0-\tfrac{11}{1152}\)

\Needspace{14\baselineskip}
Beispiel: Nr. 65. H.\ . 4 Proz.\quad \(m=53.784\)
\begin{center}
\fontsize{8}{13}\selectfont\setlength{\tabcolsep}{2pt}
\begin{tabular}{ll}
\(\log\varphi(m-\tfrac{13}{24})\ .\ .\ .\ .\ .\)&1.05488\\
\(\text{comp.}\log\rho^{\frac{1}{4}}\ .\ .\ .\ .\ .\ .\)&0.00426\\
\(\log g(m+\tfrac{11}{49})^{\frac{13}{28}}\ .\ .\ .\ .\)&0.00440\\\cline{2-2}
&1.06354\\
\(P=\)&11.5755
\end{tabular}\qquad
\begin{tabular}{ll}
\(\log\varphi m\ .\ .\ .\ .\)&1.04830\\
\(\log\rho^{\frac{7}{24}}\ .\ .\)&\(-0.00497\)\\
\(B\varphi m\ .\ .\ .\ .\)&11.0492\\
\(\tfrac{13}{24}\rho^{\frac{3}{4}}\ .\ .\ .\ .\ .\)&0.5260\\\cline{2-2}
&11.5752
\end{tabular}\par\medskip
Verbesserung des früheren Wertes \(P^0\)\par\smallskip
\(P^0-\tfrac{1}{99}\tfrac{105}{104}P^0-\tfrac{11}{1152}\)\par\smallskip
\begin{tabular}{r}
\(\tfrac{13}{24}+\varphi m=11.7179\)\\
\(-0.1314\ 5\)\\
\(-0.0095\ 5\)\\\hline
11.5769
\end{tabular}
\end{center}

\Needspace{5\baselineskip}
\pdfbookmark[2]{Bestehende Ehe.}{witwen-stehende-ehe}
\begin{center}\textit{Bestehende Ehe.}\end{center}
\noindent\hspace*{1.2em} Alter des Mannes \(M\), der Frau \(m\) zum Zeitpunkt des 1. Okt. In der Sterblichkeitstafel: lebende Männer des Alters \(x=Fx\), Wert der Witwenpension \(=R-Q\).

[Nach den Statuten beginnt die Pension mit dem Ablauf des Gnadenquartals, wird in halbjährlichen, am 1. April und am 1. Oktober jedes Jahres fälligen Raten ausgezahlt und erlischt mit dem Ende des Sterbemonats. Für den Zeitwert \(R-Q\) einer gegebenenfalls eintretenden Witwenpension, die jährlich die Geldeinheit beträgt, gilt demnach
\[
\begin{aligned}
(R-Q)\mathord{.}fm\mathord{.}Fm&=\tfrac{1}{12}\rho\{Fm-F(m+\tfrac{1}{4})\}\\
&\qquad\{f(m+\tfrac{6}{12})+f(m+\tfrac{7}{12})+f(m+\tfrac{8}{12})\}\\
&\quad+\tfrac{1}{12}\rho\{Fm-F(m+\tfrac{2}{4})\}\\
&\qquad\{f(m+\tfrac{9}{12})+f(m+\tfrac{10}{12})+f(m+\tfrac{11}{12})\}\\
&\quad+\tfrac{1}{12}\rho^{\frac{3}{2}}\{Fm-F(m+\tfrac{3}{4})\}\\
&\qquad\{f(m+\tfrac{12}{12})+f(m+\tfrac{13}{12})+f(m+\tfrac{14}{12})\}\\
&\quad+\tfrac{1}{12}\rho^{\frac{3}{2}}\{Fm-F(m+\tfrac{4}{4})\}\\
&\qquad\{f(m+\tfrac{15}{12})+f(m+\tfrac{16}{12})+f(m+\tfrac{17}{12})\}\\
&\quad+\text{usw.}
\end{aligned}
\]
\noindent oder:]
\[
\begin{aligned}
R\mathord{.}fm&=\tfrac{1}{2}\rho f(m+\tfrac{17}{24})+\tfrac{1}{2}\rho^{\frac{3}{2}}f(m+\tfrac{29}{24})+\tfrac{1}{2}\rho\rho f(m+\tfrac{41}{24})+.\ .\\
&=\rho^{\frac{5}{4}}f(m+\tfrac{23}{24})+\rho^{\frac{9}{4}}f(m+\tfrac{47}{24})+.\ .
\end{aligned}
\]
\[
R=\tfrac{1}{24}\rho^{\frac{5}{4}}+\rho^{\frac{7}{24}}\varphi m
\]
\[
\begin{aligned}
Q\mathord{.}fm\mathord{.}Fm&=\rho^{\frac{5}{4}}f(m+\tfrac{23}{24})\mathord{.}F(M+\tfrac{5}{8})\\
&\quad+\rho^{\frac{9}{4}}f(m+\tfrac{47}{24})F(M+\tfrac{13}{8})+.\ .
\end{aligned}
\]
\SourceLeaf{192}
\Needspace{5\baselineskip}
\pdfbookmark[2]{Wert der noch zu zahlenden Beiträge S.}{witwen-beitraege}
\begin{center}\textit{Wert der noch zu zahlenden Beiträge \(S\).}\end{center}
\noindent\hspace*{1.2em} [Regulativ vom 11. Oktober 1833 und 24. November 1846. Jeder Teilnehmer an der Witwenkasse hat den Beitrag jährlich nachschüssig am 17. September an den Rechnungsführer zu entrichten; diese Beiträge werden von Michaelis zu Michaelis gerechnet. — Stirbt ein Mitglied der Witwenkasse in der ersten Hälfte des Beitragsjahres, also in den Monaten von Oktober bis einschließlich März, so brauchen die Erben den Beitrag für das betreffende Jahr nicht mehr einzuzahlen; stirbt ein Mitglied dagegen in der zweiten Hälfte des Jahres, so haben die Erben am 17. September des Sterbejahres noch den vollen Beitrag zu entrichten. — Die Kündigung durch einen Teilnehmer muss durch eine schriftliche Erklärung vor dem 17. September des von Michaelis zu Michaelis laufenden Beitrittsjahres erfolgen. Der an diesem Tag fällige Jahresbeitrag muss jedoch noch einmal voll bezahlt werden. Wer diese Frist nicht einhält, muss noch für das ganze folgende Beitrittsjahr zahlen, bleibt dann aber bis dahin auch Mitglied.

Für den Zeitwert \(S\) der jährlich als Beitrag zu zahlenden Geldeinheit erhält man daher:]
\[
S\mathord{.}fm\mathord{.}Fm=\rho\mathord{.}fm\mathord{.}F(M+\tfrac{1}{2})+\rho\rho\mathord{.}f(m+1)\mathord{.}F(M+\tfrac{3}{2})+.\ .
\]
\noindent Bezeichnet man also
\[
\psi(m,M)\mathord{.}fm\mathord{.}Fm=\rho f(m+1)F(m+1)+\rho\rho f(m+2)F(m+2)+.\ .\ .
\]
\noindent so ist
\[
\begin{aligned}
Q&=\frac{f(m-\tfrac{1}{24})}{fm}\mathord{.}\frac{F(M-\tfrac{3}{8})}{FM}\mathord{.}\rho^{\frac{1}{4}}\psi(m-\tfrac{1}{24},M-\tfrac{3}{8})\\[5pt]
S&=\frac{f(m-1)}{fm}\mathord{.}\frac{F(M-\tfrac{1}{2})}{FM}\mathord{.}\psi(m-1,M-\tfrac{1}{2})
\end{aligned}
\]
\noindent Schreibt man außerdem
\[
\begin{aligned}
\frac{f(m-1)}{fm}&=gm\\[5pt]
\frac{F(M-1)}{FM}&=GM
\end{aligned}
\]
\noindent so ergibt sich
\[
\begin{aligned}
R&=\tfrac{1}{24}\rho^{\frac{5}{4}}+\rho^{\frac{7}{24}}\varphi m\\
Q&=g(m+\tfrac{23}{48})^{\frac{1}{24}}\mathord{.}G(M+\tfrac{5}{16})^{\frac{3}{8}}\mathord{.}\rho^{\frac{1}{4}}\mathord{.}\psi(m-\tfrac{1}{24},M-\tfrac{3}{8})\\
S&=gm\mathord{.}G(M+\tfrac{1}{4})^{\frac{1}{2}}\mathord{.}\psi(m-1,M-\tfrac{1}{2})
\end{aligned}
\]
\SourceLeaf{193}
\Needspace{17\baselineskip}
Beispiel: Nr. 138 Conradi.\quad \(M=65.02,\quad m=46.08\)
\begin{center}
\fontsize{9}{14}\selectfont\setlength{\tabcolsep}{5pt}
\begin{tabular}{lrr}
&4 Proz.&\(3\tfrac{1}{2}\) Proz.\\
\(\log\varphi m\ .\ .\ .\ .\)&1.12875&1.15233\\
\(\tfrac{7}{24}\log\rho\ .\ .\ .\ .\)&\(-497\)&\(-436\)\\
&13.2979&14.0594\\
\(\tfrac{1}{24}\rho^{\frac{5}{4}}\ .\ .\ .\ .\)&0.0397&399\\\cline{2-3}
\(R=\)&13.3376&14.0993\\\cline{2-3}
für \(Q,\ \log\psi\ .\ .\ .\)&0.80912&0.82183\\
\(\log gG\ .\ .\)&\(+897\)&\(+897\)\\
\(\tfrac{1}{4}\log\rho\ .\ .\ .\)&\(-426\)&\(-374\)\\
\(Q=\)&6.5137&6.7150\\
\(R-Q=\)&6.8239&7.3843\\\hline
für \(S,\ \log\psi\ .\ .\ .\)&0.81325&0.82606\\
\(\log gG\ .\ .\)&\(+1744\)&\(+1744\)\\
\(S=\)&6.7716&6.9743\\\hline
\end{tabular}
\end{center}
\[
\begin{aligned}
\log\rho&=-0.0170333393\quad\text{für 4 Proz.}\\
\log\rho&=-0.0149403498\quad\text{für }3\tfrac{1}{2}\text{ Proz.}
\end{aligned}
\]
\begin{center}
\fontsize{9}{14}\selectfont\setlength{\tabcolsep}{4pt}
\begin{tabular}{c|c}
für \(Q\)&für \(S\)\\
\(\psi(46.04;64.65)\ .\ .\ 18.61\)&\(\psi(45.08;64.52)\ .\ .\ 19.44\)\\
18.61&19.44\\
\begin{tabular}{r||r|r}64&0.82079&0.83374\\65&0.80284&0.81540\end{tabular}&\begin{tabular}{r|r}0.82251&0.83552\\0.80470&0.81732\end{tabular}\\
\(g(46.56)\ .\ .\ .\ 25\)&\(g(46.08)\ .\ .\ .\ 588\)\\
\(G(65.33)\ .\ .\ .\ 872\)&\(G(65.27)\ .\ .\ .\ 1156\)
\end{tabular}
\end{center}

\noindent\hspace*{1.2em} Bei der Aufstellung dieser Tafeln hat \textsc{Gauss} seine fünfstellige Tafel zur Berechnung des Logarithmus der Summe von Größen benutzt, die nicht selbst, sondern nur durch ihre Logarithmen gegeben sind. Deshalb hat er bei jedem Zinssatz und jedem Altersunterschied mit der Bestimmung der Rentenwerte für die höchsten Lebensjahre begonnen. Die einzelnen Zahlenangaben können aus diesem Grund um einzelne Einheiten der fünften Dezimalstelle der Logarithmen ungenau sein, auch abgesehen von der Verbesserung, die die Sterblichkeitstafel seitdem bereits durch erweiterte Beobachtungen der Preußischen Witwenkasse erfahren hat.

Einige Rechenfehler, auf die ich durch die Bildung der Quotienten zwischen den \(3\tfrac{1}{2}\)- und 4-prozentigen Rentenwerten und der Unterschiede aufeinanderfolgender Quotienten aufmerksam geworden bin, habe ich beim Abdruck berichtigt. In Einheiten der fünften Dezimalstelle des Logarithmus betragen diese Fehler an den Stellen, an denen sie entstehen: \(+20\) für das 77. Lebensjahr der Frau in ihren Leibrentenwerten bei 4 Prozent, ferner für das 71. und 93. Lebensjahr des Mannes und die jeweiligen Altersunterschiede der Frau von \(+1\) und 0 Jahren bei \(3\tfrac{1}{2}\) Prozent; \(+10\) für das 76. Lebensjahr des Mannes in seinem Leibrentenwert bei \(3\tfrac{1}{2}\) Prozent und ebenso viel für das 79. Lebensjahr des Mannes und den Altersunterschied der Frau von \(-9\) Jahren bei 4 Prozent; \(-10\) für das 28., 68., 91., 94. und 94. Lebensjahr des Mannes und die jeweiligen Altersunterschiede von \(-8\), \(-13\), \(-11\), \(-2\), und \(-14\) Jahren bei den jeweiligen Prozentsätzen \(4.4.3\tfrac{1}{2}.4\) und \(3\tfrac{1}{2}\). Diese und einige andere kleinere Rechenfehler haben die Bestimmung der Rentenwerte für die unmittelbar jüngeren Altersjahre etwas beeinflusst. Dieser Einfluss verschwindet allmählich und im äußersten Fall erst am Ende des zweiten Jahrzehnts. Die Zahlenwerte der einfachen Leibrenten sind aus den Logarithmen abgeleitet und mussten nach deren angegebener Berichtigung hier ebenfalls entsprechend geändert werden.

Die bei der Anwendung der Tafeln zu verwendenden Formeln und die zugehörigen Rechenbeispiele sind den verstreuten Notizen auf einzelnen Blättern entnommen und hier durch einige Einschübe erläutert
\begin{flushright}\textsc{Schering.}\end{flushright}
\begin{center}\rule{.3\linewidth}{.3pt}\end{center}
